% Direct-authority diplomatic French for printed p.177.
\SourcePage{177}

\LayerTarget{tohoku:corollary:3.8.1:4:p2}
\noindent\textit{vide $(U_1,\ldots\mathpunct{.}U_k)$ d'ouverts de $\mathbf U$, leur
intersection $U$ satisfait à $\check H^i(U,F)=\{0\}$ pour $i>0$. Alors on a
aussi $H^i(U,F)=\{0\}$, et pour toute partie ouverte $V$ de $X$,
l'homomorphisme naturel $\check H^i(V,F)\to H^i(V,F)$ est un
isomorphisme.}

\LayerTarget{tohoku:proof:3.8.1:corollary-4:fixed-cover}
Il suffit de prouver $H^i(U,F)=\{0\}$, car $V$ admet des recouvrements
arbitrairement fins $\mathbf R$ par des ouverts de $\mathbf U$, et on conclut
du corollaire 1 que pour un tel $\mathbf R$ l'homomorphisme
$\check H^i(R,F)\to H^i(V,F)$ est un isomorphisme, ce qui
prouve en même temps ($\mathbf R$ étant arbitrairement fin) que
$\check H^i(V,F)\to H^i(V,F)$ est un isomorphisme.
\LayerTarget{tohoku:proof:3.8.1:corollary-4:induction}
--- Pour prouver $H^i(U,F)=\{0\}$, on prouve par récurrence sur $n$ que
$H^i(U,F)=\{0\}$ pour $0<i\leqq n$ et tout $U$ de la forme indiquée. C'est
trivial si $n=0$, supposons $n\geqq 1$ et l'assertion démontrée pour
$n'=n-1$. Il y a des recouvrements arbitrairement fins $\mathbf R$ de $U$ par
des ouverts de $\mathbf U$, pour un tel $\mathbf R$ on a
\LayerTarget{tohoku:eq:corollary-4-cover-complex-vanishing}
$C(\mathbf R,H^q(F))=0$ pour $0<q<n$
d'après l'hypothèse de récurrence, à fortiori
\LayerTarget{tohoku:eq:corollary-4-cover-cohomology-vanishing}
$H^p(R,H^q(F))=0$ pour de tels $q$,
d'où
\LayerTarget{tohoku:eq:corollary-4-cech-limit-vanishing}
$\check H^p(U,H^q(F))=0$ pour de tels,
ce qui, en vertu de la remarque précédent le corollaire 4, implique que
\LayerTarget{tohoku:eq:corollary-4-induction-conclusion}
$H^n(U,F)=\check H^q(U,F) $
qui est nul.

\LayerTarget{tohoku:application:affine-zariski-coherent-sheaf}
Le corollaire 4 s'applique par exemple au cas oû $X$ est une variété
algébrique muni de sa topologie de Zariski, $\mathbf U$ l'ensemble des ouverts
affines dans $X$, $F$ un faisceau algébrique cohérent sur $X$ [15]. En effet,
d'après [15], les ouverts affines forment une base de la topologie de $X$, et
l'intersection de deux ouverts affines est un ouvert affin si $U$ est un
ouvert affine, on a $\check H^i(U,F)=0$ pour $i>0$. On a donc
\LayerTarget{tohoku:eq:affine-cech-derived-comparison}
$H^i(X,F)=\check H^i(X,F)$,
de plus le corollaire 1 ci-dessus montre que l'on peut calculer les $H^i(X,F)$
à l'aide d'un recouvrement arbitrairement choisi de $X$ par des ouverts
affines.

\LayerTarget{tohoku:remark:closed-cover-leray-spectral-sequences}
\textsc{Remarque.} Il y a d'autres cas que celui du théorème 3.8.1 oû
la suite spectrale de Leray est valable. Le plus connu est celui d'un
\textit{recouvrement localement fini de $X$ supposé paracompact par des
ensembles fermés} (c'est le cas envisagé par Leray) ; il se traite le plus
simplement comme ci-dessus pour les recouvrements ouverts, grâce au fait que
la restriction d'un faisceau mou à une partie fermé est encore un faisceau mou
(remplaçant prop. 3.1.3). Un autre cas, dégagé par Godement par une méthode
différente, est celui d'un \textit{recouvrement fini de $X$ par des ensembles
fermés} (sans hypothèse de paracompacité). Quand on est simultanément dans ces
deux cas, les deux suites spectrales obtenues coïncident, fort heureusement.

\LayerTarget{tohoku:example:cech-H2-comparison-failure:p1}
\textsc{Un exemple.} Pour terminer ce N\textsuperscript{o}, nous
allons indiquer un exemple simple oû le monomorphisme
$\check H^2(X,F)\to H^2(X,F)$ n'est pas un isomorphisme, et oû on
a même $\check H^2(X,F)=0$, $H^2(X,F)\neq 0$. Comme on tire du corollaire 2 au
théorème 3.8.1 une suite exacte
\LayerTarget{tohoku:eq:cech-H2-comparison-exact-sequence}
\[
0\to\check H^2(X,F)\to H^2(X,F)
\to\check H^1(X,H^1(F))\to 0
\]
\LayerTarget{tohoku:criterion:cech-H2-example-reduction:p1}
il suffit \textit{d'indiquer un cas où $H^2(X,F)\neq 0$ et
$H^2(X,F)\to\check H^1(X,H^1(F))$ est un}

\TohokuContinuation{178}
